% !TEX root = ../eigenfaces_algorithm.tex
%=====================================================
\section{识别：在 $k$ 维坐标里找最近邻}

\subsection{流程}

\begin{enumerate}
  \item \textbf{建画廊}：280 张训练图各自投影，得到坐标矩阵 $\mat{C}\in\R^{280\times k}$（每行一个 $\mathbf{c}_j$）及标签 $y_j$（人的编号）。
  \item \textbf{投测试图}：$\mathbf{c}_t=\mat{V}_k\mT(\mathbf{a}_t-\bar{\mathbf{a}})$。
  \item \textbf{找最近邻}：算 280 个距离 $d_j=\norm{\mathbf{c}_t-\mathbf{c}_j}_2$，取最小者的标签作为预测（平票取索引最小者）。
  \item \textbf{算准确率}：120 张测试图判对的比例。对 $k\in\{5,10,20,50,100\}$ 制表并画曲线。
\end{enumerate}

\begin{intuition}
为什么坐标距离能代替“长相相似度”？因为 $\mat{V}_k$ 的列是正交基，
内积、距离这些几何量在基变换下保持不变（正交变换不改变长度与夹角）。
$\norm{\mathbf{c}_t-\mathbf{c}_j}$ 本来就是“两张脸在特征脸子空间里的真实距离”，
只是把 4096 维里无信息的维度统统删掉了。同一个人的三张测试脸，
减去均值后剩下的“差异模式”与训练脸高度相似，因此在坐标空间里自然聚成一团。
\end{intuition}

\subsection{可选变体：白化距离}

经典 Eigenfaces 文献常用“白化”后的距离：
\[
d_j^2=\sum_{i\le k}\frac{(c_{t,i}-c_{j,i})^2}{\lambda_i} ,
\]
即先除掉每个方向的天然幅度（1 号特征脸系数天然比 50 号大几个数量级，不除会让头部方向垄断距离）。
它通常改变识别结果，\textbf{只能作为并列变体报告，不能与主结果混为一谈}。

\subsection{为什么 $k$ 不是越大越好（一条值得写进报告的解释）}

正交分解给出精确的恒等式：
\[
\norm{\mathbf{a}_t-\mathbf{a}_j}^2
=\norm{\mathbf{c}_t-\mathbf{c}_j}^2+\norm{\mathbf{r}_t-\mathbf{r}_j}^2 ,
\]
而 $k$ 维坐标距离\textbf{恰好扔掉了残差项} $\norm{\mathbf{r}_t-\mathbf{r}_j}^2$。
$k$ 小：残差大、且携带真实的判别信息（人脸轮廓差异还躺在被扔掉的方向里），
丢掉它是欠拟合；$k$ 大：残差项变小，但噪声维度的坐标进入距离并开始主导排序，
是过拟合。所以准确率随 $k$ 一般先升后平（或缓降），最佳 $k$ 常落在解释方差 $90\sim95\%$ 附近。
注意这是实验观察，报告中应表述为“实验中观察到”，而非理论必然。
